\documentclass{beamer}					% Document class

\usepackage[english]{babel}				% Set language
\usepackage[utf8x]{inputenc}			% Set encoding

\mode<presentation>						% Set options
{
  \usetheme{default}					% Set theme
  \usecolortheme{default} 				% Set colors
  \usefonttheme{default}  				% Set font theme
  \setbeamertemplate{caption}[numbered]	% Set caption to be numbered
}

% Uncomment this to have the outline at the beginning of each section highlighted.
%\AtBeginSection[]
%{
%  \begin{frame}{Outline}
%    \tableofcontents[currentsection]
%  \end{frame}
%}

\usepackage{graphicx}					% For including figures
\usepackage{booktabs}					% For table rules
\usepackage{hyperref}					% For cross-referencing
\usepackage{subcaption}  % 記得加在 preamble
\usepackage{multicol}

\title{Introduction to R}	% Presentation title
\author{Rex Shao-Yu Wang}								% Presentation author
\institute{Department of Agronomy, National Taiwan University}					% Author affiliation
\date{\today}									% Today's date	

\begin{document}

% Title page
% This page includes the informations defined earlier including title, author/s, affiliation/s and the date
\begin{frame}
  \titlepage
\end{frame}

% Outline
% This page includes the outline (Table of content) of the presentation. All sections and subsections will appear in the outline by default.
\begin{frame}{Outline}
  \tableofcontents
\end{frame}

% The following is the most frequently used slide types in beamer
% The slide structure is as follows:
%
%\begin{frame}{<slide-title>}
%	<content>
%\end{frame}

\section{Introduction}
\begin{frame}{About Me}
    \begin{itemize}
        \item Rex Wang 
        \item First-year master's student in Agronomy Biostatistics
        \item I will not be in Steven's lecture
        \item Feel free to email me if you have any questions  
        \item Email: R14621207@ntu.edu.tw
        \item Feel free to interrupt me if you have any questions
    \end{itemize}    
\end{frame}

\begin{frame}{Lab Schedule}

\begin{columns}[T]

% Left column
\column{0.35\textwidth}
\textbf{Time \& Location}

\vspace{0.5em}
\begin{itemize}
    \item Monday 13:20--14:10
    \item Boya 409
    \item One week behind theory class
\end{itemize}

% Right column
\column{0.7\textwidth}
\textbf{Weekly Topics}

\vspace{0.5em}

\small
\renewcommand{\arraystretch}{1.0}
\begin{tabular}{lll}
\toprule
\textbf{Week} & \textbf{Date} & \textbf{Topic} \\
\midrule
1  & 2/23 & No class \\
2  & 3/2  & Introduction to R \\
3  & 3/9  & Descriptive Statistics \\
4  & 3/16 & Visualization of Data \\
5  & 3/23 & Discrete Random Variables \\
6  & 3/30 & Continuous Random Variables \\
7  & 4/6  & No class (Holiday) \\
8  & 4/13 & Sampling Distribution \\
9  & 4/20 & Estimation \\
10 & 4/27 & Hypothesis Testing I \\
11 & 5/4  & Hypothesis Testing II \\
12 & 5/11 & Hypothesis Testing III \\
13 & 5/18 & ANOVA \\
14 & 5/25 & Correlation \& Simple Regression \\
15 & 6/1  & Categorical Data Analysis \\
16 & 6/8  & No class (Final Exam) \\
\bottomrule
\end{tabular}
\end{columns}
\end{frame}

\begin{frame}{Rules of Lab Session}

\begin{columns}[T]

% Left column
\column{0.65\textwidth}

\begin{itemize}
    \item No food or drink allowed!
    
    \item Lab Report: 9\% of total grade
    \begin{itemize}
        \item 8 reports this semester
        \item Available on NTU COOL
        \item Deadline: Following Monday at 23:59
        \item Plagiarism $\rightarrow$ 0\% for all involved
    \end{itemize}
    \item It should be done during lab session
\end{itemize}

% Right column
\column{0.35\textwidth}
\vspace{-30pt}

\includegraphics[width=\textwidth]{f9eje-srjou.jpg}

\end{columns}

\end{frame}


\section{R and R Studio}

\begin{frame}{R and RStudio}

\begin{columns}[T]

% Left column
\column{0.7\textwidth}

\small

\textbf{R} is a programming language used for statistical computing and data visualization. 

\vspace{0.5em}

\textbf{RStudio} is an integrated development environment (IDE) that provides a user-friendly interface for working with R.
\begin{itemize}
    \item Free
    \item Continuous improvement (R-4.5.3 now)
    \item Low barrier to coding
    \item Available at: https://posit.co/download/rstudio-desktop/
\end{itemize}

% Right column
\column{0.35\textwidth}
\vspace{0pt}

\centering

\includegraphics[width=0.7\linewidth]{R_logo.svg.png}

\vspace{1em}


\includegraphics[width=0.9\linewidth]{rstudio.png}

\end{columns}

\end{frame}

\begin{frame}{Installation}
Click on the link: https://posit.co/download/rstudio-desktop/ \\
Download R \textbf{FIRST} before RStudio!
\begin{figure}
    \centering
    \includegraphics[width=1\linewidth]{截圖 2026-02-27 上午11.12.15.png}
\end{figure}
\end{frame}

\begin{frame}{Download R}
\begin{figure}
    \centering
    \includegraphics[width=1\linewidth]{截圖 2026-02-27 上午11.11.18.png}
\end{figure}
\end{frame}

\begin{frame}{Download R}
\begin{figure}
    \centering
    \includegraphics[width=1\linewidth]{截圖 2026-02-27 上午11.15.29.png}
\end{figure}
\end{frame}

\begin{frame}{Download R}

\begin{columns}[T]

\column{0.5\textwidth}
\centering
\includegraphics[width=1.1\linewidth]{截圖 2026-02-27 上午11.32.17.png}

\column{0.5\textwidth}
\centering
\includegraphics[width=1.1\linewidth]{截圖 2026-02-27 上午11.32.54.png}

\end{columns}

\end{frame}

\begin{frame}{Download R}

\begin{columns}[T]

\column{0.5\textwidth}
\centering
\includegraphics[width=1.1\linewidth]{截圖 2026-02-27 上午11.33.30.png}

\column{0.5\textwidth}
\centering
\includegraphics[width=1.1\linewidth]{截圖 2026-02-27 上午11.33.59.png}

\end{columns}

\end{frame}

\begin{frame}{Download R}

\begin{columns}[T]

\column{0.5\textwidth}
\centering
\includegraphics[width=1.1\linewidth]{截圖 2026-02-27 上午11.34.35.png}

\column{0.5\textwidth}
\centering
\includegraphics[width=1.1\linewidth]{截圖 2026-02-27 上午11.35.45.png}

\end{columns}

\end{frame}

\begin{frame}{Download RStudio}
\begin{figure}
    \centering
    \includegraphics[width=1\linewidth]{截圖 2026-02-27 上午11.57.11.png}
\end{figure}
    
\end{frame}

\begin{frame}{Download RStudio}
Check the application to see whether it was successfully installed
\begin{columns}[T]

\column{0.5\textwidth}
\centering
\includegraphics[width=1.1\linewidth]{截圖 2026-02-27 中午12.01.17.png}

\column{0.5\textwidth}
\centering

\includegraphics[width=1.1\linewidth]{截圖 2026-02-27 中午12.02.51.png}

\end{columns}
\end{frame}

\begin{frame}{RStudio Interface}
\begin{figure}
    \centering
    \includegraphics[width=1.05\linewidth]{截圖 2026-02-27 中午12.30.58.png}
\end{figure}
\end{frame}

\begin{frame}{How to create a new R Script}
File $\rightarrow$ New file $\rightarrow$ R Script
\begin{figure}
    \centering
    \includegraphics[width=1\linewidth]{截圖 2026-02-27 下午3.40.46.png}
\end{figure}
\end{frame}

\begin{frame}{Change the Appearance}
Tools $\rightarrow$ Global Option $\rightarrow$ Appearance
\begin{columns}[T]

\column{0.5\textwidth}
\centering
\includegraphics[width=1.1\linewidth]{截圖 2026-02-27 下午3.47.31.png}

\column{0.5\textwidth}
\centering
\includegraphics[width=1.1\linewidth]{截圖 2026-02-27 下午3.47.00.png}

\end{columns}
\end{frame}

\section{Basic}
\begin{frame}{The Basics: Calculations}
- R will do math for you\\
- How to run a command: \\
\begin{itemize}
    \item Type Type command into the console and type \fbox{\texttt{Enter}}
    \item Type the command in the script and type 
    \begin{itemize}
        \item \fbox{\texttt{Ctrl}} \fbox{\texttt{Enter}} (Windows)
        \item \fbox{\texttt{Command}} \fbox{\texttt{Enter}} (macOS)
    \end{itemize}
\end{itemize}
\end{frame}

\begin{frame}{Basic Arithmetic Operations}
\begin{columns}[T]

\column{0.6\textwidth}
The basic arithmetic operators are:
\begin{itemize}
    \item + (addition)
    \item - (subtraction)
    \item * (multiplication)
    \item / (division)
    \item \texttt{\^} (exponentiation)
    \item \%/\% (quotient)
    \item \%\% (modulus)
\end{itemize}

\column{0.5\textwidth}
\centering
\includegraphics[width=0.65\linewidth]{截圖 2026-02-27 下午4.06.52.png}

\end{columns}
\end{frame}

\begin{frame}{Basic Arithmetic Function}
\begin{columns}
\column{0.6\textwidth}

\textbf{Functions} accept \textbf{parameters} (in
parentheses) and return a value. \\
\vspace{15 pt}
The basic arithmetic functions are:
\begin{itemize}
    \item log(x): logarithms base e of x
    \item exp(x): Exponential of x
    \item abs(x): absolute value of x
    \item sqrt(x): square root of x
\end{itemize}

\column{0.45\textwidth}
\begin{figure}
    \centering
    \includegraphics[width=0.9\linewidth]{截圖 2026-02-27 下午4.23.55.png}
\end{figure}
\end{columns}
    
\end{frame}

\begin{frame}{The Basic: Variables}
\begin{columns}
\column{0.75\textwidth}
\begin{itemize}
    \item Variables are named containers for data.
    \item The assignment operator in R is: \texttt{<-} or \texttt{=}
    \item Variable names can start with a letter
    \begin{itemize}
        \item Just not a number by itself.
        \item Examples: result, x1
    \end{itemize}
    \item R is case sensitive (i.e. Result \neq  \text{result})
\end{itemize}

\column{0.45\textwidth}
\begin{figure}
    \centering
    \includegraphics[width=1\linewidth]{截圖 2026-02-27 下午4.55.17.png}
\end{figure}
\end{columns}
\end{frame}

\section{Practice}
\begin{frame}{Practice }

Use R to calculate the following. \\
Try assigning values to variables first, then perform the calculation

\vspace{0.5em}

\begin{enumerate}

\item 
\[
\frac{46\times628}{71}
\]

\item 
\[
\ 7 \times (2^8 +56)
\]

\item 
\[
\frac{log(|3-5|+1)+\sqrt{3^2+5^2}-e^{3-5}}{(3^3+5^3+1)}
\]

\end{enumerate}

\end{frame}



%%%%%%%%%%%%%%%%%%%%%%%%%%%%%%%%%%%%%%%%%%%%%%%%%%%%%%%%%%



\end{document}
